% =====================================================================
% III-A: cấp độ 1. Ablation ở cuối (III-A5); kết quả test ở Mục IV-B.
\subsection{Cấp độ 1: xếp lại bằng scene graph trên VG $\cap$ COCO}
\label{sec:method-l1}

\subsubsection{Dựng đồ thị}\label{sec:l1-graphs}
\textbf{Đồ thị ảnh.} Mỗi ảnh có một scene graph riêng, lấy từ \texttt{objects.json} và \texttt{relationships.json} của
Visual Genome~\cite{krishna2017vg} (VG). Nhãn giữ dạng mở, chỉ chuẩn hoá alias, không lọc về VG150 (bước kiểm tra độ phủ cho
thấy lọc VG150 làm mất các nhãn hay gặp như \emph{water}, \emph{grass}, \emph{field}). Quan hệ có hướng; các quan hệ trùng
cả ba thành phần được gộp. Không có cạnh giữa các ảnh, nên đồ thị không chứa cạnh suy ra từ thống kê toàn tập. Mỗi ảnh
train và val có trung bình 26,7 node và 18,0 cạnh; 3,0\% ảnh không có cạnh nào.

\textbf{Đồ thị câu.} Mô hình cấp độ~1 tách caption bằng SceneGraphParser~\cite{schuster2015sgparser} (gói
\texttt{sng\_parser}): trung bình 3,06 vật thể và 1,65 quan hệ mỗi câu, và 14,2\% câu val không tách được quan hệ nào. Mô
hình chính dùng FACTUAL-T5 large~\cite{li2023factual}, một parser seq2seq đã huấn luyện sẵn để sinh scene graph từ câu.
Parser chỉ chạy suy luận, không tinh chỉnh. Trên val, FACTUAL giảm tỉ lệ node là đại từ từ 3,9\% xuống 0,1\%, và tăng tỉ lệ
câu có ít nhất một quan hệ từ 85,8\% lên 97,3\%.

\subsubsection{Bộ xếp lại GATv2}\label{sec:l1-model}
\textbf{Mã hoá đồ thị.} Đồ thị câu và đồ thị ảnh dùng chung một bộ mã hoá. Vector khởi đầu của mỗi node là vector CLIP
text của nhãn (câu mẫu ``a photo of a \{nhãn\}''), chiếu từ 512 xuống 256 chiều. Mỗi cạnh mang vector của nhãn quan hệ.
Hai lớp GATv2~\cite{brody2022gatv2} (4 head) cho mỗi node gom thông tin từ hàng xóm theo trọng số attention học được. Mỗi
bộ ba (chủ thể, quan hệ, đối tượng) đi qua một MLP. Vector cuối của mỗi phần là $\mathrm{L2}(f + U g)$, với $f$ là vector
ban đầu, $g$ là đầu ra GAT và $U$ khởi tạo bằng 0.

\textbf{Hai kênh khớp.} Kênh vật thể chỉ so vật thể với vật thể; kênh bộ ba chỉ so bộ ba với bộ ba. Với phần $i$ của câu
và phần $j$ cùng loại của ảnh, $\mathit{sim}_{ij}$ là cosine giữa hai vector. Mỗi phần của câu được so với mọi phần cùng
loại của ảnh qua attention:
\begin{equation}
a_{ij} = \operatorname{softmax}_j\!\left(\mathit{sim}_{ij}/\tau\right),\qquad s_i = \textstyle\sum_j a_{ij}\,\mathit{sim}_{ij}.
\end{equation}
Trọng số của từng phần trong câu là
$w_i = \operatorname{softmax}_i(u^\top p_i)$, với $p_i$ là vector của phần $i$. Điểm của kênh là
$S = \sum_i w_i s_i$.

\textbf{Trộn điểm.} Ba trọng số kênh tính riêng cho từng câu từ vector CLIP của câu:
$[a, b, c] = \operatorname{softmax}(W q + \mathit{bias})$. Điểm đồ thị là
$G = z\big((b\,z(S_\text{obj}) + c\,z(S_\text{tri}))/(b+c)\big)$, và điểm cuối là $a\,z(\text{CLIP}) + (1-a)\,G$. Đây
đúng là công thức của đề với $\alpha = a$, chỉ khác ở chỗ $\alpha$ thay đổi theo câu. Mọi tham số mới khởi tạo sao cho mô
hình chưa huấn luyện cho đúng phép thử không huấn luyện đã đo trước ($U=0$, $u=0$, $W=0$,
$\mathit{bias} = \log(0{,}60;\,0{,}28;\,0{,}12)$, $\tau = 0{,}05$). Mô hình có 1\,054\,214 tham số huấn luyện được; CLIP
không nằm trong số này (chi tiết: \file{docs/level1-model.md}).

\textbf{Huấn luyện.} Loss là cross-entropy trên 50 ứng viên của điểm cuối, chia cho nhiệt độ học được. AdamW, lr
$10^{-3}$, weight decay $10^{-4}$, batch 256, dropout 0,1. Epoch đầu khoá $W$ và $\mathit{bias}$ để phần đồ thị học trước.
Dừng sớm khi R@1 val không tăng sau 2 epoch, tối đa 10 epoch, giữ checkpoint tốt nhất theo val. Ứng viên của câu train lấy
trong một nhóm ảnh train cố định (Mục~\ref{sec:leakage}). Mỗi cấu hình chạy 3 seed (0, 1, 2).

\subsubsection{Mô hình chính}\label{sec:method-main}
Mô hình chính (\texttt{countonly\_factual\_r3}) giữ nguyên khung trên và thay bốn chỗ. Mỗi thay đổi được giữ lại vì làm R@1
val tăng qua 3 seed (Bảng~\ref{tab:l1-chain}).
\begin{enumerate}
\item \textbf{Đồ thị câu từ FACTUAL-T5} thay cho \texttt{sng\_parser}.
\item \textbf{Trộn điểm thô:} điểm cuối $= a\cdot\cos + b\cdot S_\text{obj} + c\cdot S_\text{tri}$, bỏ z-score. Ứng viên
không có phần nào để so ở một kênh nhận điểm trung bình của các ứng viên có phần đó. Nhờ bỏ chuẩn hoá, điểm cuối tách
\emph{chính xác} thành tổng các phần, nên lời giải thích không cần xấp xỉ.
\item \textbf{Đồ thị ảnh sạch và liền.} Hai node cùng loại (cùng synset hoặc hypernym trong WordNet~\cite{miller1995wordnet})
có IoU $\ge 0{,}7$ được gộp làm một. Mỗi cạnh có hộp ở hai đầu được gán 6 số hình học (độ lệch tâm, tỉ lệ kích thước, IoU,
tỉ lệ diện tích). Mỗi node nhận thêm cạnh từ 5 node gần nhất, và một node tổng nối tới mọi node. Cạnh thêm này chỉ truyền tin
trong GAT, không được tính là bộ ba. Sau bước này 100\% đồ thị ảnh val liền một khối (cấp độ~1: 4,7\%).
\item \textbf{Số lượng.} Mỗi node ảnh biết có bao nhiêu vật cùng loại trong ảnh (nhóm 1 / 2 / 3--4 / $\ge5$). Luật viết tay
trên câu gán mỗi vật thể một lớp số lượng (số nêu rõ, ``a pair of'', danh từ số nhiều, \dots; còn lại là ``không rõ''). Một
bảng học được $g$ ($7\times4$, khởi tạo 0, hàng ``không rõ'' cố định bằng 0) được cộng vào $\mathit{sim}_{ij}$ của kênh vật
thể.
\end{enumerate}

\subsubsection{Lời giải thích sinh từ cơ chế xếp hạng}\label{sec:method-expl}
Nhóm không chạy bước tìm đường đi nào sau khi đã có kết quả; mọi số trong lời giải thích được ghi lại trong cùng lần tính
điểm. Với mô hình chính, vì trộn điểm thô, điểm cuối tách chính xác thành
$a\cdot\cos + \sum_i b\,w_i \sum_j a_{ij}\mathit{sim}_{ij} + \sum_i c\,w_i \sum_j a_{ij}\mathit{sim}_{ij}$. Bảng giải thích
ghi với mỗi phần $i$ của câu: $w_i$, phần $j$ của ảnh có attention lớn nhất, độ khớp, phần trăm của điểm cuối và, với vật
thể có số lượng, giá trị $g$ đã cộng. Một cột nữa ghi giá trị của dòng đó trừ trung bình của chính dòng đó trên 50 ứng
viên. Tổng cột này đúng bằng điểm cuối trừ điểm trung bình, nên cột này cho biết phần đóng góp của từng dòng vào chênh lệch
giữa ảnh này và các ảnh khác. Đồ thị con câu $\leftrightarrow$ ảnh hiển thị trong demo (Mục~\ref{sec:demo}) là chính các cặp
$(i, j)$ này.

\subsubsection{Ablation}\label{sec:l1-ablation}
Mọi số trong mục này là R@1, trung bình $\pm$ độ lệch chuẩn qua 3 seed; khi so hai cấu hình, nhóm dùng kiểm định Welch trên
3 seed mỗi bên.

\begin{table}[t]
\centering
\caption{Từng bước từ mô hình cấp~1 tới mô hình chính, R@1 và MRR trên val. $p$: Welch so với bước trước.}
\label{tab:l1-chain}
\footnotesize
\setlength{\tabcolsep}{2.5pt}
\begin{tabular}{p{3.4cm}ccc}
\toprule
Bước (đổi một thứ so với bước trước) & R@1 & MRR & Chênh R@1 ($p$) \\
\midrule
Mô hình cấp 1 & \pmstd{46,48}{0,24} & 58,85 & \\
+ Đồ thị câu FACTUAL-T5 & \pmstd{47,23}{0,36} & 59,49 & +0,76 ($\approx$0,05) \\
+ Trộn điểm thô & \pmstd{47,48}{0,10} & 59,72 & +0,25 (n.s.) \\
+ Đồ thị ảnh sạch và liền\textsuperscript{*} & \pmstd{48,38}{0,17} & 60,51 & +0,90 (0,003) \\
+ Số lượng và bảng $g$ & \textbf{\pmstd{48,91}{0,21}} & \textbf{60,86} & +0,53 (0,032) \\
\bottomrule
\multicolumn{4}{p{0.95\columnwidth}}{\scriptsize \textsuperscript{*}Gộp bốn thay đổi (gộp trùng, hình học cạnh, cạnh
kNN, node tổng) theo quyết định của nhóm; ablation bỏ cạnh ở Bảng~\ref{tab:l1-ablation-main} đo phần đóng góp của cạnh.}
\end{tabular}
\end{table}

\textbf{Từng bước xây mô hình chính} (Bảng~\ref{tab:l1-chain}). Mỗi bước đổi một thứ và được giữ lại khi R@1 val tăng qua 3
seed. Mức tăng của bước trộn điểm thô không có ý nghĩa thống kê. Bước này vẫn được giữ vì nó cho phép tách chính xác điểm
cuối thành từng phần (Mục~\ref{sec:method-expl}) và có độ lệch giữa các seed nhỏ nhất. Các hướng thử trên val nhưng không
giữ (so khớp bộ ba tách thành phần, kênh mức cả đồ thị, gom vật cùng loại thành một node nhóm, \dots) ghi ở
\file{experiments/improve/README.md}. Ví dụ, gom vật cùng loại làm R@1 giảm 0,63 điểm ($p = 0{,}019$), nên mô hình chính chỉ
thêm số lượng vào từng node. Thêm thuộc tính (màu, chất liệu) vào phép thử không huấn luyện chỉ tăng R@1 val 0,36 điểm,
dưới ngưỡng 0,5 điểm nhóm đặt trước, nên nhóm không thêm thuộc tính vào mô hình
(\file{experiments/checks/attribute-probe.REPORT.md}).

\textbf{Cải thiện có đến từ cấu trúc quan hệ không.} Đề yêu cầu bỏ cạnh hoặc thay bằng cạnh ngẫu nhiên để kiểm tra điều
này. Nhóm làm hai loại ablation. Loại thứ nhất huấn luyện lại từ đầu, cùng công thức, cùng 3 seed
(Bảng~\ref{tab:l1-ablation} cho mô hình cấp~1, dòng ``bỏ mọi cạnh'' của Bảng~\ref{tab:l1-ablation-main} cho mô hình chính).
Loại thứ hai giữ nguyên checkpoint và chỉ phá đồ thị lúc chấm.

\begin{table}[t]
\centering
\caption{Ablation của mô hình cấp~1, mỗi dòng huấn luyện lại từ đầu. R@1.}
\label{tab:l1-ablation}
\footnotesize
\setlength{\tabcolsep}{3pt}
\begin{tabular}{p{4.3cm}cc}
\toprule
Cấu hình & Val & Test \\
\midrule
Mô hình cấp 1 & \pmstd{46,48}{0,24} & \pmstd{45,77}{0,23} \\
$a,b,c$ cố định cho mọi câu & \pmstd{46,13}{0,15} & \pmstd{45,13}{0,05} \\
Bỏ GAT (không lan truyền) & \pmstd{46,71}{0,07} & \pmstd{45,95}{0,20} \\
Tắt kênh bộ ba (GAT vẫn dùng cạnh) & \pmstd{45,99}{0,05} & \pmstd{45,07}{0,22} \\
Bỏ hẳn quan hệ & \pmstd{45,98}{0,10} & \pmstd{45,26}{0,14} \\
Câu không qua bộ mã hoá đồ thị & \pmstd{46,92}{0,07} & \pmstd{45,61}{0,29} \\
Nối lại cạnh ngẫu nhiên (giữ bậc) & \pmstd{46,24}{0,24} & \pmstd{45,50}{0,40} \\
Loss LambdaRank thay cross-entropy & \pmstd{45,18}{0,34} & \pmstd{44,28}{0,30} \\
\bottomrule
\end{tabular}
\end{table}

\begin{table}[t]
\centering
\caption{Ablation và đối chứng của mô hình chính, chỉ trên val. R@1 (chênh so với mô hình đủ).}
\label{tab:l1-ablation-main}
\footnotesize
\setlength{\tabcolsep}{3pt}
\begin{tabular}{p{3.5cm}cc}
\toprule
 & Mô hình chính & Mô hình cấp 1 \\
\midrule
Đủ đồ thị & \pmstd{48,91}{0,21} & \pmstd{46,48}{0,24} \\
\multicolumn{3}{l}{\emph{Huấn luyện lại}} \\
Bỏ mọi cạnh & $-$2,39 ($p$ = 0,0024) & $-$0,50 ($p$ = 0,055) \\
\multicolumn{3}{l}{\emph{Chỉ phá lúc chấm}} \\
Đồ thị của một ảnh khác & 21,91 ($-$27,00) & 26,33 ($-$20,15) \\
Nối lại quan hệ có nhãn ngẫu nhiên & 48,32 ($-$0,59) & 46,29 ($-$0,19) \\
Bỏ bảng $g$ & $-$0,35 & --- \\
\bottomrule
\end{tabular}
\end{table}

Bảng~\ref{tab:l1-ablation} và~\ref{tab:l1-ablation-main} cho bốn kết luận.
\begin{itemize}
\item \textbf{Mô hình dựa vào đồ thị của đúng ảnh.} Khi mỗi ảnh mang scene graph của một ảnh khác, R@1 val của mô hình chính
giảm từ 48,91 xuống 21,91, thấp hơn CLIP thuần (39,13).
\item \textbf{Ở mô hình cấp~1, cạnh đóng góp ít.} Bỏ hẳn quan hệ làm R@1 giảm khoảng 0,5 điểm; bỏ GAT không làm giảm kết quả.
\item \textbf{Ở mô hình chính, cạnh đóng góp 2,39 điểm.} Huấn luyện lại sau khi bỏ mọi cạnh làm R@1 val giảm 2,39 điểm
($p = 0{,}0024$), tương ứng 24\% mức chênh giữa mô hình chính và CLIP trên val (2,39 / 9,78); ở mô hình cấp~1 tỉ lệ này là
7\%. Khi bỏ cạnh, mô hình chính còn 46,52, ngang mô hình cấp~1 (46,48). Xáo các quan hệ có nhãn lúc chấm làm R@1 giảm 0,59
điểm. Phần đóng góp riêng của từng loại cạnh (quan hệ có nhãn, kNN, node tổng) chưa được tách.
\item \textbf{Phần mất khi bỏ cạnh tập trung ở câu có quan hệ và câu nhiều vật thể}; câu chỉ có vật thể gần như không đổi
(\file{experiments/improve/no-edges/REPORT.md}).
\end{itemize}
